\subsection*{Binding a request to the order of pictured clauses}
Four questions on the same dictionary card can contain exactly the same glyphs yet ask for four different answers. One exchange assigns the expressions for already knowing and wanting to know to the original meaning and the pronunciation. Another exchange determines which historical form should be shown. A reader invariant to glyph permutation can recover at most one of the four complete answers. The original writer's features were identical across all four alternatives on 1,135 of 1,136 development groups.

A 51,776-parameter model read the pictured question in order and wrote 4,542 of 4,544 exact answers on held-out families. It recovered all four answers on 1,135 groups. A matched reader whose added encoder received the same glyphs in sorted order reached 1,136 answers and zero complete groups (Fig.~\ref{fig:question-order}). Both preserved all 1,136 original simple answers. Replacing the status expressions with reserved wording and reversing the outer clause order gave 4,540 exact answers and 1,134 complete groups. Each status still followed the field it modified.

In a subsequent test, moving the same status expressions before their respective fields preserved every glyph count and the requested answer. Accuracy fell to 356 of 4,544, with zero complete groups. This positional failure places a limit on the interpretation of the high development score. The reader learned a local grammatical convention that determines both the lexical fact and the historical ink to write. The transfer measured here does not establish general interpretation of Chinese requests.
